\begin{table}[t]
\centering
\caption{First Stage: Inspector Leniency Predicts Appeal Outcomes}
\label{tab:first_stage}
\begin{tabular}{lccc}
\hline\hline
 & Coefficient & $N$ & Fixed Effects \\
\hline
\multicolumn{4}{l}{\textit{Panel A: First Stage (Dep. Var.: Appeal Allowed)}} \\[3pt]
Leniency (cell-specific) & -0.0538 & 860 & LPA$\times$Type + Year$\times$Type \\
 & (0.0226) & & \\
Leniency (simple FEs) & -0.0414 & 857 & LPA + Type + Year \\
 & (0.0194) & & \\
First-stage $F$ & 5.6 / 4.6 & & \\[6pt]
\multicolumn{4}{l}{\textit{Panel B: Balance Tests (Dep. Var. in Left Column)}} \\[3pt]
Filing lag (days) & -1.7523 & 857 & \\
 & (3.0783) & & \\
Householder type (=1) & -0.0103 & 857 & \\
 & (0.0205) & & \\
\hline
\multicolumn{4}{p{0.85\textwidth}}{\footnotesize \textit{Notes:} Panel~A reports the first stage of the inspector leniency IV. The dependent variable is an indicator for the appeal being allowed. Leniency is the leave-one-out inspector allow rate within case-type$\times$year cells, standardized to mean zero and unit variance. Panel~B reports balance tests: leniency should not predict case characteristics within cells. Standard errors clustered at LPA level in parentheses.} \\
\end{tabular}
\end{table}
